% fig_vl22_boltzmann.tex — Was der Faktor e^{-beta E_n} mit einem festen
% Niveauschema macht, wenn man die Temperatur dreht: von der Spitze im
% Grundzustand bis in die Gleichverteilung.
%
% WHY THIS FIGURE.  The lecture derives rho = e^{-beta H}/Z and then
% names e^{-E_j/kT} the Boltzmann factors — but the whole physical
% content of that formula is a *shape*, and a shape is what prose is
% worst at.  Two things the text states but never shows:
%   * the same fixed ladder of levels is occupied completely differently
%     at different T, and the two ends of the temperature scale are the
%     two states this course already knows — the pure ground state
%     (S = 0) and the maximally mixed state (S = k log d);
%   * the boxed entropy formula S = k beta <H> + k log(Z) is a curve.
%     Read at T -> 0 it is the third law that closes this lecture:
%     S -> 0.  Read at T -> infinity it saturates at k log d.
% Panel (a) is the microscopic side (which level carries what weight),
% panel (b) the macroscopic side (what a thermometer and a calorimeter
% would report).  That is the contrast the two panels exist for.
%
% GEOMETRY FROM THE FORMULA.  Nothing is drawn by eye; every bar, curve
% and asymptote is evaluated from the closed forms of this lecture on
% one concrete spectrum: the equally spaced ladder of Vorlesung 12,
% truncated to d = 6 levels and counted from the ground state,
%
%       E_n = n Delta,   n = 0,1,...,5,
%
% (equal spacing, because then Z is a finite geometric series and the
% populations are an exact geometric sequence — the constant ratio
% p_{n+1}/p_n = e^{-beta Delta} is what makes the fanning-out of the
% bars readable).  With x = e^{-Delta/kT} = e^{-beta Delta},
%
%       Z    = sum_{n=0}^{5} x^n = (1 - x^6)/(1 - x),
%       p_n  = x^n / Z = e^{-n Delta/kT} (1-x)/(1-x^6),
%       <H>/Delta = (sum_n n x^n)/Z,
%       S/k  = beta <H> + log(Z)          [the boxed formula of the lecture]
%
% and both limits are exact:
%       T -> 0        p_0 -> 1,   <H> -> 0,        S -> 0,
%       T -> infinity p_n -> 1/d, <H> -> (d-1)Delta/2, S -> k log d.
%
% The three temperatures drawn are kT/Delta = 0.3, 1, 3.  Their exact
% populations (the bar lengths) are
%
%   kT/Delta   x        Z        p_0     p_1     p_2     p_3    p_4    p_5
%     0.3   0.035674  1.036994  0.9643  0.0344  0.0012  4.4e-5  2e-6   6e-8
%     1     0.367879  1.578055  0.6337  0.2331  0.0858  0.0315  0.0116 0.0043
%     3     0.716531  3.050300  0.3278  0.2349  0.1683  0.1206  0.0864 0.0619
%   (1/d = 0.1667 is the T -> infinity value of every one of them.)
%
% and the points marked on the curves of (b) are
%
%   kT/Delta      0.3       1        3      -> limit
%   <H>/Delta   0.0370   0.5671   1.5886      2.5
%   S/k         0.1596   1.0233   1.6448      log 6 = 1.7918
%
% The bars are drawn by pgfplots from these formulas (xcomb at samples
% n = 0..5), not typed in; the numbers above are recorded so the plate
% can be checked without re-running LaTeX.
%
% DESIGN DECISIONS
%   * colour is temperature and nothing else, on the ORDERED ramp the
%     templates keep for ordered parameters (tpclevels):
%     spacecadet -> cgblue -> birrenorange, cold to hot.  It was chosen
%     over the obvious blue-orange-red because those three are NOT
%     monotone in luminance: printed grey, cgblue (0.40) and birrenred
%     (0.34) come out the same, and the legend of (a) would then have
%     two identical rows.  This ramp runs 0.15 -> 0.40 -> 0.54, so cold
%     stays darkest in grey as well.  The same three colours mark the
%     same three temperatures on the curves of (b), which is how one
%     reads off what the hot distribution costs in entropy;
%   * the two curves of (b) are therefore NOT on that ramp: they are
%     birrenred, a colour the plate uses nowhere else, and they share
%     it, because they are two quantities of one and the same state —
%     what tells them apart is the dash pattern and the label written
%     on each.  The temperature dots on them carry a white ring so they
%     read against the curve in colour and in grey alike, and the two
%     dots belonging to one temperature are joined by a short stroke in
%     that temperature's colour — without it the pair at kT = 3 Delta,
%     0.06 apart, fuses into a single blob;
%   * within every level of (a) the bars run cold-top to hot-bottom,
%     always in that order, so the grouping can be read without the
%     legend;
%   * the energy levels are neutral grey rules (levellinegrey): they are
%     the reference scheme, the populations are the data;
%   * grey dotted verticals (classicallim) mark the two limits of the
%     probability axis, p = 1/d and p = 1, and are labelled by the
%     limit that produces them, not by their value.  They are reference
%     lines, hence grey, hence the same style as the asymptotes of (b);
%   * a thin exponential envelope p(E) = e^{-beta E}/Z accompanies each
%     bar group, over exactly the range of the spectrum.  It is the
%     continuous function the levels sample, and it carries the same
%     vertical offset as the bars it belongs to, so it
%     passes exactly through their tips (the offset is legibility only —
%     the true energies are the grey rules);
%   * probability runs horizontally and energy vertically, so the level
%     scheme stands the way every other level diagram in these notes
%     does (Vorlesungen 11, 12, 16).

\begin{center}
\begin{tikzpicture}

% ─────────────────────────────────────────────────────────────
%  The formulas of the lecture, once, as macros.  #1 is kT/Delta.
%  \boltzP is p_n with n as the plot variable x (so it doubles as the
%  continuous envelope p(E) with E = x Delta).
% ─────────────────────────────────────────────────────────────
\def\boltzP#1{ exp(-x/#1)*(1-exp(-1/#1))/(1-exp(-6/#1)) }
% Z, <H>/Delta and S/k with kT/Delta as the plot variable x.
\def\partZ{ ((1-exp(-6/x))/(1-exp(-1/x))) }
\def\meanU{ ((exp(-1/x)+2*exp(-2/x)+3*exp(-3/x)+4*exp(-4/x)+5*exp(-5/x))/\partZ ) }
\def\entrS{ ((1/x)*\meanU + ln(\partZ)) }

% ─────────────────────────────────────────────────────────────
%  (a)  The populations on a fixed ladder, at three temperatures
% ─────────────────────────────────────────────────────────────
\begin{axis}[
    name=pop,
    tpcplot,
    width=7.2cm, height=5.2cm,
    grid=none,
    clip=false,
    xmin=0, xmax=1.13, ymin=-0.55, ymax=6.25,
    xtick={0,0.1666667,0.5,1},
    xticklabels={$0$,$1/d$,$0.5$,$1$},
    ytick={0,1,2,3,4,5},
    xlabel={$p_n$},
    ylabel={$E_n/\Delta$},
    every axis y label/.style={at={(axis description cs:-0.11,0.5)},
                               rotate=90, anchor=south, font=\small},
    title={(a) \tpcifenglish{populations}{Besetzungen}
              $p_n\propto e^{-\beta E_n}$},
    legend style={at={(0.99,0.99)}, anchor=north east, font=\scriptsize},
    legend cell align=left,
]

  % The level scheme E_n = n Delta: reference rules, not data.
  \draw[levellinegrey] (axis cs:0,0) -- (axis cs:1.10,0);
  \draw[levellinegrey] (axis cs:0,1) -- (axis cs:1.10,1);
  \draw[levellinegrey] (axis cs:0,2) -- (axis cs:1.10,2);
  \draw[levellinegrey] (axis cs:0,3) -- (axis cs:1.10,3);
  \draw[levellinegrey] (axis cs:0,4) -- (axis cs:1.10,4);
  \draw[levellinegrey] (axis cs:0,5) -- (axis cs:1.10,5);

  % The two limits of the probability axis.
  \draw[classicallim] (axis cs:0.1666667,-0.45) -- (axis cs:0.1666667,5.65);
  \draw[classicallim] (axis cs:1,-0.45)         -- (axis cs:1,5.65);

  % Envelopes p(E) = e^{-beta E}/Z, each on its own group's offset.
  \addplot[spacecadet,   line width=0.6pt, forget plot,
           domain=0:5, samples=140] ({\boltzP{0.3}}, {x+0.26});
  \addplot[cgblue,       line width=0.6pt, forget plot,
           domain=0:5, samples=140] ({\boltzP{1}},   {x});
  \addplot[birrenorange, line width=0.6pt, forget plot,
           domain=0:5, samples=140] ({\boltzP{3}},   {x-0.26});

  % The populations themselves, sampled at n = 0,...,5.
  \addplot[xcomb, mark=none, spacecadet,   line width=3.4pt,
           domain=0:5, samples=6] ({\boltzP{0.3}}, {x+0.26});
  \addlegendentry{$kT=0.3\,\Delta$}
  \addplot[xcomb, mark=none, cgblue,       line width=3.4pt,
           domain=0:5, samples=6] ({\boltzP{1}},   {x});
  \addlegendentry{$kT=\Delta$}
  \addplot[xcomb, mark=none, birrenorange, line width=3.4pt,
           domain=0:5, samples=6] ({\boltzP{3}},   {x-0.26});
  \addlegendentry{$kT=3\,\Delta$}

  % What the spectrum is.
  \node[lbl, font=\scriptsize, text=birrengray, anchor=west]
    at (axis cs:0.02,5.95) {$E_n=n\Delta$, $d=6$};

  % The limits, written under the value they force.
  \node[lbl, font=\scriptsize, text=birrengray, anchor=north]
    at (axis cs:0.1666667,-1.22) {$T\to\infty$};
  \node[lbl, font=\scriptsize, text=birrengray, anchor=north]
    at (axis cs:1,-1.22) {$T\to0$};
\end{axis}

% ─────────────────────────────────────────────────────────────
%  (b)  The same state seen from outside: <H> and S against T
% ─────────────────────────────────────────────────────────────
\begin{axis}[
    tpcplot,
    at={(pop.south east)}, anchor=south west, xshift=14mm,
    width=5.6cm, height=5.2cm,
    xmin=0, xmax=6.4, ymin=0, ymax=2.85,
    xtick={0,1,2,3,4,5,6}, ytick={0,1,2},
    xlabel={$kT/\Delta$},
    title={(b) \tpcifenglish{energy and entropy}{Energie und Entropie}},
]
  % Saturation values, the T -> infinity ends of both curves.
  \addplot[levellinegrey, domain=0:6.4, samples=2] {2.5};
  \addplot[levellinegrey, domain=0:6.4, samples=2] {1.7917595};
  \node[lbl, font=\scriptsize, text=birrengray, anchor=south west]
    at (axis cs:0.15,2.5) {$(d{-}1)/2$};
  \node[lbl, font=\scriptsize, text=birrengray, anchor=south west]
    at (axis cs:0.15,1.7917595) {$\log d$};

  % At each of the three temperatures of (a), the two curves are read at
  % once; a short stroke in that temperature's colour joins the pair, so
  % it reads as one temperature with two values.  Needed most at
  % kT = 3 Delta, where the readings differ by 0.06 and two bare dots
  % fuse into a single blob — the stroke turns that near-coincidence
  % into what it is: the two curves about to cross.  Kept short on
  % purpose; a full-height vertical would run through the annotations.
  \draw[spacecadet,   line width=0.6pt] (axis cs:0.3,0.00) -- (axis cs:0.3,0.20);
  \draw[cgblue,       line width=0.6pt] (axis cs:1,0.53)   -- (axis cs:1,1.06);
  \draw[birrenorange, line width=0.6pt] (axis cs:3,1.56)   -- (axis cs:3,1.67);

  % <H>/Delta and S/k.  Neutral: these are quantities, not temperatures.
  \addplot[birrenred, line width=1.0pt, domain=0.05:6.4, samples=220]
    {\meanU};
  \addplot[birrenred, line width=1.0pt, densely dashed,
           domain=0.05:6.4, samples=220] {\entrS};

  \node[lbl, anchor=south east] at (axis cs:6.35,2.06)
    {$\langle\widehat{H}\rangle/\Delta$};
  \node[lbl, anchor=north east] at (axis cs:6.35,1.70) {$S/k$};

  % The readings themselves, at the ends of those strokes.
  \node[circle, fill=spacecadet,   inner sep=0.9pt, draw=white, line width=0.4pt] at (axis cs:0.3,0.0370) {};
  \node[circle, fill=spacecadet,   inner sep=0.9pt, draw=white, line width=0.4pt] at (axis cs:0.3,0.1596) {};
  \node[circle, fill=cgblue,       inner sep=0.9pt, draw=white, line width=0.4pt] at (axis cs:1,0.5671) {};
  \node[circle, fill=cgblue,       inner sep=0.9pt, draw=white, line width=0.4pt] at (axis cs:1,1.0233) {};
  \node[circle, fill=birrenorange, inner sep=0.9pt, draw=white, line width=0.4pt] at (axis cs:3,1.5886) {};
  \node[circle, fill=birrenorange, inner sep=0.9pt, draw=white, line width=0.4pt] at (axis cs:3,1.6448) {};

  \node[lbl, font=\scriptsize, text=birrengray, anchor=east, align=right]
    at (axis cs:6.30,0.74)
    {\tpcifenglish{marks: the three}{Marken: die drei}\\
     \tpcifenglish{temperatures of (a)}{Temperaturen aus (a)}};

  % The third law, where this lecture ends.
  \node[lbl, font=\scriptsize, text=birrengray, anchor=west]
    at (axis cs:0.55,0.13) {$T\to0\ \Rightarrow\ S\to0$};
\end{axis}

\end{tikzpicture}
\end{center}
